\section{Methodology}
\label{sec:methodology}

\subsection{Problem Setup and Notation}
We consider generative modeling on the binary hypercube \(\{0,1\}^n\) with \(n=12\), and evaluate models on discovery of unseen states under controlled holdout protocols.
For a model distribution \(q\), a holdout set \(H \subseteq \{0,1\}^n\), and sampling budget \(Q\), discovery is measured by expected unique-state recovery
\begin{equation}
R(Q; q, H) \;=\; \frac{1}{|H|}\sum_{x\in H}\left(1-(1-q(x))^Q\right).
\end{equation}
The main scalar endpoint is
\begin{equation}
Q_{80}(q,H) \;=\; \min\{Q\in\mathbb{N}: R(Q; q,H)\ge 0.8\},
\end{equation}
with search cap \(Q_{\max}=2\times 10^5\) and \(Q_{80}=\infty\) if the threshold is not reached.
We also report holdout mass \(q(H)=\sum_{x\in H}q(x)\), enrichment \(q(H)/q_{\mathrm{unif}}(H)\), and fixed-budget recovery values (e.g., \(R(1000)\), \(R(10000)\)).

\subsection{Target Distribution Family and Region Definitions}
All experiments use the \texttt{paper\_even} target family implemented in \texttt{iqp\_generative/core.py}. Let
\begin{equation}
\Omega_{\mathrm{even}} \;=\; \{x\in\{0,1\}^n : \mathrm{parity}(x)=0\}.
\end{equation}
Define a discrete score
\begin{equation}
s(x) \;=\; 1 + \mathrm{longest\_zero\_run\_between\_ones}(x),
\end{equation}
and target distribution
\begin{equation}
p^*(x) \;\propto\; \exp(\beta s(x))\mathbf{1}[x\in\Omega_{\mathrm{even}}].
\end{equation}
We study \(\beta\)-sweeps (typically \(0.6,\ldots,1.4\)) and fixed-\(\beta\) diagnostics.

Two region notions are used:
\begin{enumerate}
\item Global region \(G_{\mathrm{global}}=\Omega_{\mathrm{even}}\).
\item High-value region \(G_{\alpha}\): top \(\alpha\) fraction by score inside \(\Omega_{\mathrm{even}}\), with default \(\alpha=0.05\).
\end{enumerate}

\subsection{Holdout Construction Protocols}
For each \((\beta,\text{seed})\), we construct holdout sets of fixed size \(|H|=20\) from a candidate region \(G\).

\paragraph{Primary protocol: Global+Smart.}
This is the default comparison protocol in the claim stack.
Candidate region is global (\(G=\Omega_{\mathrm{even}}\)), while selection is \emph{smart} via \texttt{select\_holdout\_smart}:
\begin{enumerate}
\item Apply a probability-floor relaxation \(\tau\in\{1/m,\;0.5/m,\;0.25/m,\;0\}\) until enough candidates remain (\(m\) is the holdout-construction sample scale, default \(5000\)).
\item Rank candidates by \(p^*(x)\), keep a pool (default size \(400\)).
\item Select \(H\) by greedy farthest-point sampling in Hamming distance, with tie-break by larger \(p^*(x)\).
\end{enumerate}
This yields globally scoped but non-degenerate holdouts with adequate mass and diversity.

\paragraph{High-value holdout.}
Candidate region is \(G_{\alpha}\), and the same smart-selection routine is applied.
This protocol probes discovery in concentrated high-score regions.

\paragraph{Sensitivity protocol: Global+Random.}
Candidate region is global (\(G=\Omega_{\mathrm{even}}\)), and \(H\) is sampled uniformly at random without replacement from \(G\).
This protocol is used for robustness checks and proposition-style expectation tests.

\paragraph{Train distribution under holdout masking.}
Given holdout \(H\), the train target is
\begin{equation}
p_{\mathrm{train}}(x) \;=\;
\frac{p^*(x)\mathbf{1}[x\notin H]}{1-p^*(H)}.
\end{equation}
All models are trained only from samples drawn from \(p_{\mathrm{train}}\).

\subsection{Models and Training Objectives}
We compare IQP-based and classical baselines under matched data and budget constraints.

\paragraph{IQP-QCBM.}
The primary model uses an IQP architecture with ZZ interactions on ring nearest-neighbor plus next-nearest-neighbor couplings, depth \(L=1\).
Primary loss is parity-moment MSE (\texttt{parity\_mse}); additional IQP losses (\texttt{prob\_mse}, \texttt{xent}, \texttt{mmd}) are evaluated in dedicated analyses.

\paragraph{Classical baselines.}
We include:
\begin{enumerate}
\item Ising+fields (NN+NNN), moment-matching objective.
\item Dense Ising+fields, cross-entropy objective.
\item Autoregressive Transformer, maximum likelihood objective.
\item MaxEnt parity model, parity-moment matching objective.
\end{enumerate}
In comparative claims, training data source, random seeds, and principal budget classes are matched to avoid protocol confounds.

\paragraph{Feature supervision and empirical moments.}
For parity-feature analyses, masks \(\alpha_k\in\{0,1\}^n\) are sampled i.i.d. with
\begin{equation}
p(\sigma) = \frac{1}{2}\left(1-e^{-1/(2\sigma^2)}\right),
\end{equation}
excluding \(\alpha=0\).
Given sampled masks and empirical train distribution \(\hat p_{\mathrm{train}}\), moments are
\begin{equation}
z_k = \sum_x \phi_{\alpha_k}(x)\hat p_{\mathrm{train}}(x),\qquad
\phi_{\alpha}(x)=(-1)^{\alpha\cdot x}.
\end{equation}

\subsection{Spectral Visibility and Linear Completion Analysis}
Define indicator Walsh coefficients for any set \(A\subseteq\{0,1\}^n\):
\begin{equation}
\hat 1_A(\alpha) = 2^{-n}\sum_{x\in A}\phi_{\alpha}(x).
\end{equation}
For band \(\mathcal{B}=\{\alpha_k\}_{k=1}^K\), visibility is
\begin{equation}
\mathrm{Vis}_{\mathcal{B}}(A) = \sum_{k=1}^K z_k \hat 1_A(\alpha_k).
\end{equation}
Linear reconstruction is
\begin{equation}
q_{\mathrm{lin}}(x)=2^{-n}\left(1+\sum_{k=1}^K z_k\phi_{\alpha_k}(x)\right),
\end{equation}
followed by simplex completion
\begin{equation}
\tilde q(x)\propto \max(0,q_{\mathrm{lin}}(x)).
\end{equation}
The implemented identity check is
\begin{equation}
q_{\mathrm{lin}}(H)=|H|2^{-n}+\mathrm{Vis}_{\mathcal{B}}(H),
\end{equation}
validated numerically in our spectral diagnostics.

\paragraph{Expected-scaling proposition checks.}
For random \(H\subset G\), \(|H|=h\), sampled uniformly from \(G\), we test:
\begin{align}
\mathbb{E}_H[\hat 1_H(\alpha)] &= \frac{h}{|G|}\hat 1_G(\alpha),\\
\mathbb{E}_H[\mathrm{Vis}_{\mathcal B}(H)] &= \frac{h}{|G|}\mathrm{Vis}_{\mathcal B}(G),\\
\mathbb{E}_H[q_{\mathrm{lin}}(H)] &= h2^{-n}+\frac{h}{|G|}\mathrm{Vis}_{\mathcal B}(G).
\end{align}
These are evaluated by Monte Carlo holdout sampling with fixed \(h\), across seeds and selected \(\beta\) values.

\subsection{Discovery and Fit Metrics}
\paragraph{Primary discovery metrics.}
\begin{enumerate}
\item \(Q_{80}\) (lower is better).
\item \(R(Q)\) curves over fixed \(Q\)-grid up to \(10^4\).
\item Holdout mass \(q(H)\) and enrichment \(q(H)/q_{\mathrm{unif}}(H)\).
\end{enumerate}

\paragraph{Analytic predictors and bounds.}
We report the budget-law predictor and lower bound from holdout mass:
\begin{equation}
Q_{80}^{\mathrm{pred}} \approx \frac{|H|}{q(H)}\ln 5,\qquad
Q_{80}^{\mathrm{lb}}=\frac{\ln(0.2)}{\ln(1-q(H)/|H|)}.
\end{equation}

\paragraph{Fit metrics (secondary).}
Distribution-fit diagnostics include TV and optional KL:
\begin{equation}
\mathrm{TV}(p^*,q)=\frac{1}{2}\sum_x |p^*(x)-q(x)|.
\end{equation}
Fit is treated as auxiliary to discovery.

\paragraph{Score-bucket and support-diversity metrics.}
For full-support diagnostics we use:
\begin{enumerate}
\item Score-bucket TV (\(\mathrm{TV}_{\mathrm{score}}\)).
\item BSHS(\(Q\)): bucket-weighted support-hit score.
\item Composite score: \(\mathrm{BSHS}(Q)\cdot (1-\mathrm{TV}_{\mathrm{score}})\).
\item State-level concentration and diversity metrics (dominance, \(N_{\mathrm{eff}}/|A_s|\), support-hit per score level).
\end{enumerate}

\subsection{Experimental Design and Reproducibility}
\paragraph{Data and seeds.}
Experiments are run on \texttt{paper\_even} only.
Common seed schedules are fixed (e.g., 42--46 in main sweeps), with paired seeds across model families.

\paragraph{Budget classes.}
Main comparative regimes include low-data and higher-data classes (e.g., \(m\in\{200,1000,5000\}\) in beta sweeps), with fixed holdout size (\(|H|=20\)) and pool size (400) unless overridden by script arguments.

\paragraph{Public entry points.}
Reproducible claim-level entry points are in \texttt{experiments/claims/claim01\_...\ claim11\_...}, each writing claim-scoped outputs under \texttt{outputs/claims/}.
Curated final artifacts are stored under \texttt{outputs/paper\_even\_final/}.

\subsection{Statistical Analysis}
Unless noted otherwise, statistics are aggregated over seeds for fixed \((\beta,\text{protocol},\text{model})\) tuples.
We report means and standard deviations, and for skewed/heavy-tailed discovery metrics also medians and interquartile summaries.
Pairwise model comparisons are performed on matched \((\beta,\text{seed})\) instances to preserve within-condition fairness.

For visibility-causal analyses we use rank-based tests (Spearman), with permutation-based \(p\)-values where implemented.
For proposition checks we report absolute and relative equation residuals and lhs-vs-rhs calibration plots.

\subsection{Primary vs Sensitivity Protocols}
Our primary claim-level comparisons use \textbf{Global+Smart}, because it preserves global scope while avoiding low-mass random holdout degeneracies.
Global+Random is retained as a sensitivity protocol and for expectation-level theory checks.
In particular, under global-random holdouts with very small \(p^*(H)\), uniform sampling can exceed \(p^*\) on state-coverage metrics; this behavior is expected and is explicitly separated from primary protocol conclusions.
